% Figure description placeholders — "blackboard diagrams"
% Each figure contains a detailed textual description inside a framed box.
% Replace with actual images when available: \includegraphics[width=\linewidth]{filename}

%% ===================================================================
%% 1. Ground Station Dashboard Screenshot
%% ===================================================================
\begin{figure}[htbp]
  \centering
  \fbox{\parbox{0.92\linewidth}{%
    \textit{%
      Browser window (dark theme, \texttt{http://192.168.1.3:8090}) divided into two
      main panels. \textbf{Left panel} (approximately two-thirds width): live MJPEG
      video stream from the downward-facing IMX296 camera at 640$\times$480 display
      resolution. A green bounding box is drawn around a detected rescue dummy with
      the label ``dummy 0.94'' in green text above the box. A small pink dot marks the
      detection centre. A translucent heads-up overlay in the top-left corner shows
      current telemetry: altitude (28.3\,m), ground speed (4.2\,m/s), battery voltage
      (22.4\,V), GPS fix type (3D Fix, 12 sats), and flight mode (GUIDED).
      %
      \textbf{Right panel} (one-third width, stacked vertically): \textbf{Top}---a 2D
      GPS grid (north-up) rendered on a light background with 10\,m grid squares.
      A blue arrow icon at the current drone position indicates heading. Red circles
      of varying size represent detection clusters, where circle radius encodes the
      number of detections at that location. A green diamond marks the weighted
      centroid of all confirmed-target estimates. The search polygon boundary is drawn
      as a dashed grey outline. \textbf{Bottom}---a row of five command buttons:
      \textsf{N}\,Investigate (blue), \textsf{Y}\,Confirm (green),
      \textsf{I}\,Interest (yellow), \textsf{X}\,False Positive (red),
      \textsf{L}\,Land (orange). Below the buttons, a status bar displays the current
      state machine state (e.g.\ ``CENTERING''), elapsed mission time, and number of
      detections logged.
    }%
  }}
  \caption{Web-based ground station dashboard served by \texttt{pi\_flight.py}.
    The operator views the live camera feed with detection overlays (left),
    monitors the drone's position relative to the search area and detection
    clusters (right), and issues target-classification commands via the button
    row. The entire interface runs in a standard browser with no client-side
    installation, enabling ground-station access from any device on the local
    network.}
  \label{fig:dashboard-screenshot}
\end{figure}


%% ===================================================================
%% 2. Hardware Assembly Photo
%% ===================================================================
\begin{figure}[htbp]
  \centering
  \fbox{\parbox{0.92\linewidth}{%
    \textit{%
      Oblique three-quarter photograph of the fully assembled Hexsoon EDU-450
      hexacopter on a workbench. The carbon-fibre frame has six motor arms
      radiating symmetrically, each with a brushless motor and folding propeller.
      %
      \textbf{Centre of the top plate}: the Cube Orange+ autopilot is mounted on
      a vibration-damping platform, its orange cube carrier clearly visible.
      Immediately behind it, a \textbf{Raspberry Pi~5} sits in a 3D-printed PLA
      mount secured with nylon standoffs; a short ribbon cable connects the Pi's
      CSI port downward through a slot in the top plate to the
      \textbf{IMX296 global-shutter camera} mounted on the underside, pointing
      nadir (straight down) through a rectangular cutout in the bottom plate.
      %
      A \textbf{GPS mast} (Here3+ GNSS) extends roughly 10\,cm above the frame on
      a carbon-fibre rod, positioned to minimise magnetic interference from the
      power distribution board. A \textbf{FrSky X8R receiver} is zip-tied to one
      arm with its two diversity antennas splayed at 90\textdegree.
      A \textbf{4S 5200\,mAh LiPo battery} is strapped underneath the bottom plate
      with a Velcro strap, its XT60 connector routed to the power module.
      %
      Labelled arrows (white text on semi-transparent dark background) point to
      each component: (1)~Cube Orange+, (2)~Raspberry Pi~5, (3)~IMX296 camera,
      (4)~Here3+ GPS, (5)~FrSky receiver, (6)~LiPo battery, (7)~UART cable
      (Pi $\leftrightarrow$ Cube serial connection at 921\,600\,baud),
      (8)~USB-C power for Pi.
    }%
  }}
  \caption{Hardware assembly of the SAR drone showing all major components.
    The Raspberry Pi~5 communicates with the Cube Orange+ autopilot over a
    921\,600\,baud UART link, while the IMX296 global-shutter camera feeds
    frames to the onboard YOLOv8n detector via the CSI ribbon cable.}
  \label{fig:hardware-photo}
\end{figure}


%% ===================================================================
%% 3. Detection Overlay Example
%% ===================================================================
\begin{figure}[htbp]
  \centering
  \fbox{\parbox{0.92\linewidth}{%
    \textit{%
      A single frame captured from the downward-facing camera during a search pass
      at approximately 35\,m altitude. The background is a mottled green grass field
      with scattered patches of bare earth and a gravel path crossing diagonally in
      the upper-right quadrant. Mild lens-barrel distortion is visible at the frame
      edges (corrected in software but uncorrected here for illustration).
      %
      Near the centre-left of the frame, a \textbf{rescue dummy} is visible as a
      small dark figure (roughly 40$\times$15 pixels at this altitude) lying on the
      grass. A \textbf{green bounding box} (2\,px stroke) tightly encloses the dummy.
      Above the top-left corner of the box, a label reads ``dummy 0.94'' in green
      monospace font, indicating the class name and confidence score.
      A \textbf{pink filled circle} (5\,px radius) marks the geometric centre of the
      bounding box, representing the point used for GPS back-projection.
      %
      \textbf{Top-left overlay text} (white on semi-transparent black):
      ``ALT: 34.7\,m\quad SPD: 4.1\,m/s\quad GPS: 51.42341, $-$2.67138''.
      \textbf{Bottom-left}: a horizontal \textbf{scale bar}---a white line labelled
      ``5\,m''---computed from the current altitude and calibrated field of view
      (DJI video camera HFOV = 54.4\textdegree\ for the 1456$\times$1088 frame), giving the reader an
      immediate sense of ground-sample distance. At 35\,m altitude the full frame
      covers approximately 37\,m $\times$ 28\,m on the ground.
    }%
  }}
  \caption{Example detection frame from a search pass at 35\,m altitude.
    The YOLOv8n model identifies the rescue dummy with 94\% confidence.
    The pink centre dot is back-projected through the calibrated camera model
    to produce a GPS estimate of the target location.}
  \label{fig:detection-overlay}
\end{figure}


%% ===================================================================
%% 4. Lawnmower Pattern on Satellite Image
%% ===================================================================
\begin{figure}[htbp]
  \centering
  \fbox{\parbox{0.92\linewidth}{%
    \textit{%
      Top-down satellite image of the Fenswood Farm wilderness area (University of
      Bristol flight-test site) at approximately 1\,:\,2\,000 scale. The terrain is
      a mix of open grassland, scrub, and scattered trees. A small lake is visible in
      the north-west corner.
      %
      \textbf{Green polygon outline} (solid, 2\,px): the five-vertex search area
      defined in \texttt{config.py} (\texttt{SEARCH\_AREA\_GPS}), an irregular
      pentagon covering roughly 200\,m $\times$ 150\,m of open ground.
      %
      \textbf{Blue polyline}: the computed lawnmower (boustrophedon) flight path.
      Parallel east--west legs are spaced 25.7\,m apart (80\% of the camera swath at
      30\,m altitude), with short connecting segments at alternating ends. An arrow
      on the first leg indicates the flight direction (west to east for the first
      pass). Approximately 18 waypoints are marked as small blue circles at each
      turn. The pattern begins and ends near the takeoff point to minimise
      repositioning time.
      %
      \textbf{Red polygon} (hatched fill): the SSSI (Site of Special Scientific
      Interest) no-fly zone to the south-east of the search area. The geofence
      repulsive buffer (10\,m inside the SSSI boundary) is shown as a dashed red
      line parallel to the SSSI edge.
      %
      \textbf{Yellow filled circle}: the takeoff/landing location
      (51.42341\textdegree\,N, $-$2.67145\textdegree\,W), positioned on the gravel
      access track at the western edge of the site.
      %
      A north arrow and 50\,m scale bar are placed in the bottom-right corner.
    }%
  }}
  \caption{Lawnmower search pattern overlaid on satellite imagery of the
    Fenswood Wilderness test site. The boustrophedon path (blue) covers the
    search polygon (green) with parallel legs spaced at 80\% of the camera
    footprint width, guaranteeing full visual coverage with 20\% overlap.
    The SSSI no-fly zone (red) is avoided by the geofence module.}
  \label{fig:lawnmower-satellite}
\end{figure}


%% ===================================================================
%% 5. Bench Test Setup Photo
%% ===================================================================
\begin{figure}[htbp]
  \centering
  \fbox{\parbox{0.92\linewidth}{%
    \textit{%
      Indoor workbench photograph taken from a slightly elevated angle showing the
      complete bench-test configuration.
      %
      \textbf{Centre of bench}: a Raspberry Pi~5 in its official red/white case,
      powered by a USB-C supply. A short CSI ribbon cable connects to the
      \textbf{IMX296 global-shutter camera module} resting face-up on the bench,
      pointed at a printed A3 image of the rescue dummy taped to the wall 1\,m away.
      A serial UART cable (colour-coded dupont wires: red, black, white, green) runs
      from the Pi's GPIO header to a \textbf{Cube Orange+} autopilot sitting on its
      carrier board nearby, simulating the in-flight wiring.
      %
      \textbf{Left monitor}: a laptop screen displaying \textbf{Mission Planner}
      with the HUD showing artificial horizon, compass, and telemetry readout
      (connected via TCP to mavproxy on the Pi at \texttt{tcpin:0.0.0.0:5762}).
      The Flight Data tab is active, showing ``GUIDED'' mode and 12-satellite GPS
      fix (from SITL running in the background).
      %
      \textbf{Right monitor} (or second laptop screen): a PuTTY SSH terminal
      connected to the Pi (\texttt{192.168.1.3}), showing the output of
      \texttt{python passive\_watch.py} with detection log lines scrolling:
      timestamps, confidence values, and estimated GPS coordinates.
      %
      On the bench surface, a \textbf{printed checkerboard} (A3, 14$\times$9
      squares, 28\,mm pitch) used for lens distortion calibration is propped against
      a book stand. A \textbf{tape measure} extends from the camera lens to the
      dummy printout, used for FOV calibration (92\,cm visible at 1\,m distance
      $\rightarrow$ focal length 5.46\,mm).
      %
      Cables are labelled in the photograph: (a)~CSI ribbon, (b)~UART serial,
      (c)~USB-C power, (d)~Ethernet to router.
    }%
  }}
  \caption{Bench-test setup used for integration testing before flight.
    The Raspberry Pi~5 runs the detection pipeline on live camera frames
    while communicating with the Cube Orange+ over UART. Mission Planner
    on the laptop provides telemetry monitoring via the mavproxy TCP bridge.
    The checkerboard and tape measure were used for lens and FOV calibration
    respectively.}
  \label{fig:bench-setup}
\end{figure}
